\documentclass[11pt,a4paper,sans]{moderncv}
\moderncvstyle{classic}
\moderncvcolor{blue}
\usepackage[utf8]{inputenc}
\usepackage[scale=0.85]{geometry}
\usepackage{enumitem}

% Información personal
\name{Edgar Duván}{Bernal Acero}
\title{Científico de Datos (Junior)}
\address{Madrid, Cundinamarca, Colombia}{}
\phone[mobile]{+57 313-425-1767}
\email{dvnber10@gmail.com}
\homepage{github.com/dvnber10}
\social[linkedin]{edgar-duvan-bernal-acero-43339a258}
\extrainfo{Portafolio: \url{https://portfolio-iota-seven-h96xw48djf.vercel.app/}}

\begin{document}
\makecvtitle

\section{Perfil Profesional}

\cvitem{}{\small Ingeniero en Sistemas y Computación con experiencia profesional en machine learning, pipelines de datos y modelos de IA en Grupo Consiti (proyecto StradiaLabs). Experiencia práctica en NLP aplicado a textos médicos, reconocimiento de voz en español (Whisper, Common Voice), visión por computadora y evaluación rigurosa de modelos (precisión, robustez y generalización). Sólida base en Python, Hugging Face, Librosa y bases de datos SQL/NoSQL. En vía de especialización en analítica y ciencia de datos con promedio acumulado de 4.0. Reconocido por pensamiento analítico, autonomía en investigación y foco en resultados medibles.}

\section{Logros Destacados}

\cvitem{}{\small Machine learning en producción (Grupo Consiti $\cdot$ StradiaLabs)}
\cvitem{}{\small NLP médico: 90\% de precisión en entidades y detección de negación}
\cvitem{}{\small Evaluación de modelos: precisión, robustez y generalización}
\cvitem{}{\small Especialización en analítica y ciencia de datos $\cdot$ Promedio 4.0}

\section{Educación}

\cventry{2021 – 2026 (En curso)}{Ingeniería en Sistemas y Computación}{Universidad de Cundinamarca}{Facatativá, Colombia}{}{\small
\begin{itemize}[label=$\cdot$]
    \item Noveno semestre $\cdot$ Graduación esperada: Septiembre 2026. Promedio acumulado: 4.0. Cursos: Desarrollo de Software, Inteligencia Artificial, Sistemas Operativos, Seguridad Informática. Proyecto de grado: especialización en analítica y ciencia de datos.
\end{itemize}}


\cventry{2027 (En curso)}{Especialización en Analítica y Ciencia de Datos}{Universidad de Cundinamarca}{Facatativá, Colombia}{}{\small
\begin{itemize}[label=$\cdot$]
    \item Opción de grado del pregrado $\cdot$ Graduación esperada: Marzo 2027.
\end{itemize}}



\section{Experiencia Profesional}

\cventry{Junio 2026}{Desarrollador de Inteligencia Artificial Junior}{Grupo Consiti S.A. de C.V.}{San Salvador, El Salvador (remoto)}{StradiaLabs}{\small
\begin{itemize}[label=$\cdot$]
    \item Desarrollé e implementé componentes de inteligencia artificial y machine learning para el proyecto StradiaLabs, dentro del área de Tecnología y Desarrollo.
    \item Implementé modelos y pipelines de datos orientados a proyectos de IA, entregando componentes de desarrollo dentro de los plazos acordados.
    \item Aporté soluciones técnicas orientadas a la escalabilidad del backlog del proyecto, trabajando de forma remota con autonomía y foco en resultados.
    \item Colaboré en un entorno de desarrollo distribuido bajo metodologías ágiles, con comunicación técnica efectiva en un contexto de trabajo remoto internacional.
\end{itemize}}



\section{Experiencia y Proyectos}

\cventry{2023 – 2024}{Desarrollador Back-End, Proyectos .NET Core}{Proyectos Académicos y Personales}{GitHub: Api-.net-\_y\_Mongo, MAVE}{MAVE \$\textbackslash{}cdot\$ Api-.net-\textbackslash{}\_y\textbackslash{}\_Mongo, https://github.com/dvnber10}{\small
\begin{itemize}[label=$\cdot$]
    \item Diseñé e implementé APIs RESTful en .NET Core con MongoDB, optimizando autenticación y consultas para sistemas de gestión empresarial, logrando un rendimiento 20\% superior en tiempos de respuesta.
    \item Desarrollé el backend de "MAVE", un sistema de gestión de recursos humanos, integrando endpoints seguros y configuraciones en entornos Linux/Windows.
    \item Contribuciones públicas en GitHub: 3 estrellas en MAVE y 6 repositorios activos.
\end{itemize}}


\cventry{2025}{Investigador, Modelo de Voz a Texto (ASR)}{Proyecto de Investigación Personal}{Hugging Face $\cdot$ Common Voice $\cdot$ Whisper}{ASR Español}{\small
\begin{itemize}[label=$\cdot$]
    \item Desarrollé un modelo de reconocimiento de voz en español usando Common Voice (Hugging Face) y Whisper, con preprocesamiento en Librosa para normalización de audio.
    \item Implementé enfoques híbridos de PLN para textos médicos, logrando un 90\% de precisión en reconocimiento de entidades y detección de negación.
    \item Adapté datasets para acentos colombianos, explorando VoxPopuli para robustez en entornos ruidosos.
\end{itemize}}


\cventry{2025}{Desarrollador, Sistema de Dibujo Gestual con Kinect}{Proyecto Personal}{Python $\cdot$ PyQt5 $\cdot$ Kinect $\cdot$ Tesseract OCR}{Lienzo Interactivo Kinect}{\small
\begin{itemize}[label=$\cdot$]
    \item Implementé un lienzo interactivo en Python con PyQt5 y Kinect, utilizando Tesseract OCR para extraer texto manuscrito con 85\% de precisión tras optimizar el preprocesamiento (desenfoque, dilatación).
    \item Resolví problemas de estabilidad en datos de profundidad, reduciendo interrupciones USB en un 40\% mediante ajustes en FPS y manejo de errores.
    \item Exploré integración serial con ESP32 en Arch Linux para comunicación hardware-software.
\end{itemize}}


\cventry{2024}{Automatización en Linux, Walls-Arch}{Proyecto Personal}{Arch Linux $\cdot$ feh $\cdot$ Cloudinary}{Walls-Arch, https://github.com/dvnber10/Walls-Arch}{\small
\begin{itemize}[label=$\cdot$]
    \item Creé un script en Python con feh para automatizar rotación de fondos de pantalla en Arch Linux, optimizando flujos de trabajo para desarrolladores.
    \item Demostré habilidades en scripting shell y configuración de sistemas Linux (Arch, Debian).
\end{itemize}}



\section{Proyectos Relevantes}

\cventry{}{Modelo de Voz a Texto (ASR)}{IA $\cdot$ Lenguaje}{}{}{\small
\begin{itemize}[label=$\cdot$]
    \item Reconocimiento de voz en español con Whisper y Common Voice. Logré un 90\% de precisión en reconocimiento de entidades y detección de negación en textos médicos.
    \item https://github.com/dvnber10/AppDeepLearning
\end{itemize}}


\cventry{}{ML PetMonitor}{IA $\cdot$ Machine Learning}{}{}{\small
\begin{itemize}[label=$\cdot$]
    \item API de predicción de enfermedades en mascotas mediante machine learning y procesamiento de datos del comportamiento de las mascotas.
    \item https://github.com/dvnber10/ML-PetMonitor
\end{itemize}}



\section{Habilidades Técnicas}

\cvitem{Habilidades}{\small Machine Learning • Deep Learning • Transfer Learning • Fine-tuning • NLP • ASR (Whisper) • Evaluación de modelos • Python (PyTorch, scikit-learn, Transformers, Librosa) • Ingeniería de características • Pipelines de datos • Hugging Face • MongoDB • MySQL • MS SQL Server • Git/GitHub • Docker • Experimentación}

\section{Idiomas}

\cvitem{Español}{\small Nativo}
\cvitem{Inglés}{\small B2 $\cdot$ Competente en lectura técnica, en proceso de mejora}

\section{Intereses}

\cvitem{}{\small NLP aplicado a salud e investigación, reconocimiento de voz multidioma, visión por computadora, evaluación de modelos y datos, pipelines escalables y proyectos sostenibles.}

\end{document}